\input{sol_pre.tex}
\input{q_2026-03-01_1.tex}
\begin{document}
\input{sol_head.tex}

%================= 1
\Savol{1}
$R=\underbrace{77\ldots7}_{2025}$ bölsin. Oxirgi raqami $d$ bölgan son $R+(d-7)$, \ $d=9,8,\dots,1$ (9 ta qöşiluvçi):
\[S=9R+\bigl((9+8+\dots+1)-9\cdot7\bigr)=9R+45-63=9R-18.\]
$R$ ning har bir raqami 7 $\Rightarrow$ $R=7\cdot\underbrace{11\ldots1}_{2025}$ 7 ga bölinadi. \ $S\equiv-18\equiv3\pmod7$.
\Javob{D) 3}

%================= 2
\Savol{2}
$[\lg5]=0$, \ $\{\lg0{,}1\}=\{-1\}=0$, \ $[\ln3]=1$: \quad $\{0+0-7\}=\{-7\}=0$.
\Javob{A) 0}

%================= 3
\Savol{3}
Perimetri berilgan şakllar içida eng katta yuzli — doira: $2\pi r=40\Rightarrow r=\frac{20}\pi$,
\[S=\pi r^2=\frac{400}\pi.\]
\Javob{C) $\frac{400}\pi$}

%================= 4
\Savol{4}
Tuz: $30\cdot0{,}05=1{,}5$ l. \ $\dfrac{1{,}5}{30+x}=0{,}02\Rightarrow30+x=75\Rightarrow x=45$.
\Javob{C) 45}

%================= 5
\Savol{5}
$3\cdot24^{\frac12}=3\cdot2\sqrt6=6\sqrt6=6^{\frac32}$ $\Rightarrow$ $\left(6^{\frac32}\right)^{\frac13}=\sqrt6$; \quad $2^{\frac13}\cdot54^{\frac16}=(4\cdot54)^{\frac16}=216^{\frac16}=\sqrt6$.
\[\left(\sqrt6+\sqrt6\right)^2=24.\]
\Javob{B) 24}

%================= 6
\Savol{6}
$5^7\cdot5^{\frac57}=5^{\frac{54}7}$, \ $\sqrt[6]{5^{\frac{54}7}}=5^{\frac97}$; \quad $5^{\frac97}\cdot5^{-\frac27}=5$.
\Javob{C) 5}

%================= 7
\Savol{7}
$t=\sqrt[3]2$, $t^3=2$: \ $\sqrt[3]{32}=2t^2$, $\sqrt[3]{128}=4t$. \ $(t^2+1)(-t^2+2t+1)=-t^4+2t^3+2t+1=5$ (çunki $t^4=2t$, $t^3=2$).
\[\frac5{2(t^2+1)}=\frac{-t^2+2t+1}2,\qquad\frac5{2(-t^2+2t+1)}=\frac{t^2+1}2.\]
Yiğindi: $\frac{2t+2}2=t+1$; \ $t+1-\sqrt[3]2=1$.
\Javob{B) 1}

%================= 8
\Savol{8}
$a_2=2a_1\Rightarrow d=a_1$; \ $a_{10}=10a_1=20\Rightarrow a_1=d=2$. \ $S_6-S_5=a_6=2+5\cdot2=12$.
\Javob{A) 12}

%================= 9
\Savol{9}
$b_1=a_1+d$, $b_2=a_1+8d$, $b_3=a_1+43d$; \ $b_2^2=b_1b_3$ $\Rightarrow$ $21d^2=28a_1d$ $\Rightarrow$ $a_1=\frac34d$.\\
$b_1+b_2+b_3=\frac{7+35+175}4d=\frac{217}4d=217\Rightarrow d=4$, $a_1=3$.\\
$S_n=\dfrac{6+4(n-1)}2n=n(2n+1)=820\Rightarrow n=20$.
\Javob{B) 20}

%================= 10
\Savol{10}
$1-\dfrac{x^2+y^2-4}{2xy}=\dfrac{4-(x-y)^2}{2xy}=\dfrac{(2-x+y)(2+x-y)}{2xy}$, \quad $\dfrac{\frac12+\frac1{x-y}}{\frac12-\frac1{x-y}}=\dfrac{x-y+2}{x-y-2}$.
\[\frac{(2-x+y)(x-y+2)}{2xy}\cdot\frac{x-y-2}{x-y+2}\cdot\frac{2xy}{2-x+y}=x-y-2.\]
\Javob{A) $x-y-2$}

%================= 11
\Savol{11}
$u=a+1012$, $v=b+1013$: \ $u+v=a+b+2025$.
\[(u+v)^2=4uv\iff(u-v)^2=0\iff u=v\ \Rightarrow\ a+1012=b+1013,\ \ a-b=1.\]
\Javob{A) 1}

%================= 12
\Savol{12}
$\cos2x=1-2\sin^2x$, \ $s=\sin x$: \ $\sqrt2-2\sqrt2s^2-6s-3\sqrt2\le0\iff\sqrt2s^2+3s+\sqrt2\ge0$.
\[(s+\sqrt2)\left(s+\tfrac{\sqrt2}2\right)\ge0;\quad s+\sqrt2>0\ \Rightarrow\ \sin x\ge-\tfrac{\sqrt2}2.\]
$x\in\left[-\frac\pi4+2\pi k;\ \frac{5\pi}4+2\pi k\right]$.
\Javob{C}

%================= 13
\Savol{13}
$\sin40^\circ-\cos40^\circ=\sqrt2\sin(40^\circ-45^\circ)=-\sqrt2\sin5^\circ$, \ $\sqrt2\cos85^\circ=\sqrt2\sin5^\circ$. \ Nisbat: $-1$.
\Javob{A) $-1$}

%================= 14
\Savol{14}
$\left(\frac{2025}{2026}\right)^{x^2-4}=1\Rightarrow x^2-4=0\Rightarrow x=\pm2$. \ Köpaytma: $-4$.
\Javob{C) $-4$}

%================= 15
\Savol{15}
$u=5^x$, $v=3^y$: \ $3u+2v=81$, \ $-5u+4v=-113$.\\
Birinçisini 2 ga köpaytirib ayiramiz: $11u=275\Rightarrow u=25$, $x=2$; \ $v=3$, $y=1$. \ $x+y=3$.
\Javob{D) 3}

%================= 16
\Savol{16}
$t=\sqrt{x^2-5x+4}\ge0$, \ $x^2-5x=t^2-4$: \ $2(t^2-4)+2t=4\Rightarrow t^2+t-6=0\Rightarrow t=2$.\\
$x^2-5x+4=4\Rightarrow x=0;\ 5$. \ Yiğindi: 5.
\Javob{C) 5}

%================= 17
\Savol{17}
$x\in[2;6]$. Çap tomon ösuvçi ($\sqrt{x-2}$ ösadi, $\sqrt{6-x}$ kamayadi) $\Rightarrow$ köpi bilan bitta ildiz.\\
$x=6$: $2-0=2$ — ildiz. Yiğindi: 6.
\Javob{A) 6}

%================= 18
\Savol{18}
$x^2-4x+3=(x-1)(x-3)$:
\[\frac{2x^2-14x+6-(3x-8)(x-1)}{(x-1)(x-3)}=\frac{-(x+1)(x+2)}{(x-1)(x-3)}\ge0\iff\frac{(x+1)(x+2)}{(x-1)(x-3)}\le0.\]
$x\in[-2;-1]\cup(1;3)$. \ Butunlar: $-2,-1,2$ — 3 ta.
\Javob{C) 3 ta}

%================= 19
\Savol{19}
$x\ge-6$; surat $\ge0$. \ Surat nol: $x=-6$, $x=1$. \ Maxraj manfiy: $(x-5)(x-7)<0\Rightarrow5<x<7$, $x=6$.\\
Butunlar: $-6,\ 1,\ 6$ — 3 ta.
\Javob{A) 3 ta}

%================= 20
\Savol{20}
$\frac a9+b=\frac{52}9$, \ $9a+b=-12$ $\Rightarrow$ $\frac{80}9a=-\frac{160}9$, \ $a=-2$, \ $b=6$.\\
$y=6-2\cdot3^x$, \ $3^x>0\Rightarrow E(y)=(-\infty;6)$. Teskari funksiya aniqlaniş sohasi $=E(y)$.
\Javob{A) $(-\infty;6)$}

%================= 21
\Savol{21}
$\dfrac1{x^4+1}=\dfrac12\left(\dfrac{1+\frac1{x^2}}{x^2+\frac1{x^2}}-\dfrac{1-\frac1{x^2}}{x^2+\frac1{x^2}}\right)$; \ $u=x-\frac1x$: $x^2+\frac1{x^2}=u^2+2$; \ $v=x+\frac1x$: $x^2+\frac1{x^2}=v^2-2$.
\[\frac12\int\frac{du}{u^2+2}-\frac12\int\frac{dv}{v^2-2}=\frac1{2\sqrt2}\arctg\frac{x^2-1}{\sqrt2x}+\frac1{4\sqrt2}\ln\left|\frac{v+\sqrt2}{v-\sqrt2}\right|+C.\]
$\arctg\frac{x^2-1}{\sqrt2x}$ va $\arctg\frac{\sqrt2x}{1-x^2}$ özgarmasga farq qiladi ($\arctg w+\arctg\frac1w=\pm\frac\pi2$).
\Javob{A}

%================= 22
\Savol{22}
$f'(x)=e^{x-7}(1+x-8)=e^{x-7}(x-7)=0\Rightarrow x=7$ (minimum).\\
$f(5)=-\frac3{e^2}$, \ $f(7)=-1$, \ $f(9)=e^2$. \ Eng kiçigi $-1$.
\Javob{A) $-1$}

%================= 23
\Savol{23}
Aylana $B$ dan ötib, $AC$ ga $H$ örtada urinadi $\Rightarrow$ $BH$ — diametr: $BH=6\sqrt3$, $r=3\sqrt3$, markaz $O$ — $BH$ örtasi.\\
$\angle BKH=90^\circ$ ($BH$ diametr) $\Rightarrow$ $BK=BH\cos30^\circ=9$. \ $\triangle OBK$: $OB=OK=r$, $\angle OBK=30^\circ\Rightarrow\angle BOK=120^\circ$.
\fig{\begin{tikzpicture}[scale=0.33]
\coordinate (A) at (0,0); \coordinate (C) at (12,0); \coordinate (B) at (6,10.392); \coordinate (O) at (6,5.196);
\coordinate (K) at (1.5,2.598); \coordinate (L) at (10.5,2.598);
\fill[okgreen!30] (B)--(K) arc[start angle=210,end angle=330,radius=5.196]--cycle;
\draw (A)--(B)--(C)--cycle; \draw (O) circle (5.196);
\draw[axs] (O)--(K); \draw[axs] (O)--(L); \draw[axs] (O)--(B); \draw[ax] (O)--(6,0);
\fill (O) circle (3pt);
\node[below left] at (A){$A$}; \node[below right] at (C){$C$}; \node[above] at (B){$B$}; \node[left] at (K){$K$}; \node[right] at (L){$L$};
\node[right] at (O){$O$}; \node[below] at (6,0){$H$};
\end{tikzpicture}}
Böyalgan soha $=$ $120^\circ$ li sektor $KOL$ $+$ $\triangle OBK+\triangle OBL$:
\[S=\frac{120}{360}\pi\cdot27+2\cdot\frac12\cdot27\sin120^\circ=9\pi+\frac{27\sqrt3}2=\frac{18\pi+27\sqrt3}2.\]
\Javob{D) $\frac{18\pi+27\sqrt3}2$}

%================= 24
\Savol{24}
I: $f(-x)g(-x)=-f(x)g(x)$ — toq. \ II: $\frac{f(x)}{-g(x)}$ — toq. \ III: $f(x)-g(x)$ — umuman olganda toq emas.
\Javob{B) 2 ta}

%================= 25
\Savol{25}
$AB=6a$, $BC=6b$. Böyalgan soha $=S_{BQS}-S_{BPR}$, bunda $BP=2a$, $BQ=5a$, $BR=3b$, $BS=4b$:
\[S_{BQS}=\frac{5a}{6a}\cdot\frac{4b}{6b}\cdot108=60,\qquad S_{BPR}=\frac{2a}{6a}\cdot\frac{3b}{6b}\cdot108=18,\qquad S=42.\]
\Javob{A) 42}

%================= 26
\Savol{26}
Içki aylanalar uriniş nuqtalarigaça: $DM=\frac{AD+BD-AB}2$, \ $DN=\frac{AD+DC-AC}2$.\\
$AB=AC$ bölgani uçun:
\[MN=DM-DN=\frac{BD-DC}2=\frac{14-6}2=4.\]
\Javob{B) 4}

%================= 27
\Savol{27}
$C$ dan $CE\parallel AB$ ötkazamiz: $\triangle ECD$ tomonlari $17,\ 25,\ 44-16=28$.
\fig{\begin{tikzpicture}[scale=0.14]
\coordinate (A) at (0,0); \coordinate (D) at (44,0); \coordinate (B) at (8,15); \coordinate (C) at (24,15); \coordinate (E) at (16,0);
\draw (A)--(B)--(C)--(D)--cycle; \draw[axs] (C)--(E); \draw[ax] (C)--(24,0);
\node[below left] at (A){$A$}; \node[below right] at (D){$D$}; \node[above] at (B){$B$}; \node[above] at (C){$C$}; \node[below] at (E){$E$};
\node[above] at (16,15){$16$}; \node[left] at (4,7.5){$17$}; \node[right] at (34,7.5){$25$}; \node[right,font=\scriptsize,aux] at (24,6){$h$};
\end{tikzpicture}}
Geron: $p=35$, \ $S=\sqrt{35\cdot18\cdot10\cdot7}=210$ $\Rightarrow$ $h=\dfrac{2\cdot210}{28}=15$. \quad $S_{ABCD}=\dfrac{16+44}2\cdot15=450$.
\Javob{B) 450}

%================= 28
\Savol{28}
$\dfrac{n(n-3)}2=35\Rightarrow n=10$. Muntazam önburçakka taşqi aylana çizamiz, har bir tomon $36^\circ$ yoyni tortadi.\\
Bir uçdan eng kiçik diagonal 2-qöşni uçga, eng kattasi (diametr) 5-uçga boradi. Ular orasidagi içki çizilgan burçak $3\cdot36^\circ=108^\circ$ yoyga tiraladi:
\[\varphi=\frac{108^\circ}2=54^\circ.\]
\Javob{A) $54^\circ$}

%================= 29
\Savol{29}
$\overrightarrow{OE}=\overrightarrow{OB}+\overrightarrow{BE}$, \ $\overrightarrow{OB}=\frac12\overrightarrow{DB}=\frac12\left(\overrightarrow{AB}-\overrightarrow{BC}\right)$, \ $\overrightarrow{BE}=\frac34\overrightarrow{BC}$:
\[\overrightarrow{OE}=\frac12\overrightarrow{AB}+\frac14\overrightarrow{BC}\ \Rightarrow\ a+b=\frac34.\]
\Javob{D) $\frac34$}

%================= 30
\Savol{30}
$S_{toʻla}=2S_{yon}\Rightarrow S_{asoslar}=S_{yon}$: \ $\pi(16+144)=\pi(4+12)l\Rightarrow l=10$.\\
$h=\sqrt{l^2-(12-4)^2}=6$. \quad $V=\dfrac{\pi h}3(R^2+Rr+r^2)=2\pi(144+48+16)=416\pi$.
\Javob{C) $416\pi$}

%================= 31
\Savol{31}
Kamida bittasiga: $25-10=15$. \ Ikkalasiga: $9+8-15=2$.
\Javob{A) 2}

%================= 32
\Savol{32}
«MATEMATIKA» — 10 harf, ulardan M — 2 ta. \ $P=\frac2{10}=\frac15$.
\Javob{C) $\frac15$}

%================= 33-35
\newpage
\Savol{33}
$O$ — asos markazi ($SO\perp ABC$). $ED$ — örta çiziq: $ED\parallel AC$, $ED=1$. $OM\perp ED$ bölsin.
\fig{\begin{tikzpicture}[scale=2.2]
\coordinate (A) at (0,0); \coordinate (C) at (2,0); \coordinate (B) at (1,1.732); \coordinate (O) at (1,0.577);
\coordinate (E) at (0.5,0.866); \coordinate (D) at (1.5,0.866); \coordinate (M) at (1,0.866);
\draw (A)--(B)--(C)--cycle; \draw (E)--(D);
\draw[ax] (B)--(1,0); \draw[axs] (O)--(M); \draw[ax] (O)--(D);
\fill (O) circle (0.8pt);
\node[below left] at (A){$A$}; \node[below right] at (C){$C$}; \node[above] at (B){$B$};
\node[left] at (E){$E$}; \node[right] at (D){$D$}; \node[below right] at (O){$O$}; \node[above right,font=\scriptsize] at (M){$M$};
\end{tikzpicture}}
$B$ dan $AC$ gaça $\sqrt3$, $B$ dan $ED$ gaça $\frac{\sqrt3}2$, \ $BO=\frac23\sqrt3$ $\Rightarrow$ $OM=\frac{2\sqrt3}3-\frac{\sqrt3}2=\frac{\sqrt3}6$.\\
Uç perpendikulyar teoremasi: $SM\perp ED$, \ $\angle SMO=45^\circ$ $\Rightarrow$ $SO=OM=\frac{\sqrt3}6$, \ $SM=\frac{\sqrt3}6\sqrt2=\frac{\sqrt6}6$.\\
\textbf{33.} $OD=r=\frac{\sqrt3}3$: \ $SD=\sqrt{\frac1{12}+\frac13}=\sqrt{\frac5{12}}=\frac{\sqrt{15}}6$.\\
\textbf{34.} $S_{SED}=\frac12\cdot ED\cdot SM=\frac12\cdot1\cdot\frac{\sqrt6}6=\frac{\sqrt6}{12}$.\\
\textbf{35.} $S_{AEDC}=S_{ABC}-S_{EBD}=\sqrt3-\frac{\sqrt3}4=\frac{3\sqrt3}4$; \ $V=\frac13\cdot\frac{3\sqrt3}4\cdot\frac{\sqrt3}6=\frac18$.
\Javob{33 — D; \quad 34 — F; \quad 35 — C}

%================= 36
\newpage
\Savol{36}
$a=\sqrt3$ deb, tenglamani $a$ ga nisbatan kvadrat tenglama sifatida yozamiz:
\[a^2-(2x^2+1)a+(x^4-x)=0,\qquad D=(2x^2+1)^2-4(x^4-x)=(2x+1)^2.\]
\[a=\frac{2x^2+1\pm(2x+1)}2\ \Rightarrow\ \sqrt3=x^2+x+1\quad\text{yoki}\quad\sqrt3=x^2-x.\]
1) $x^2+x+1-\sqrt3=0$: \ $D=4\sqrt3-3>0$, köpaytma $1-\sqrt3<0$ $\Rightarrow$ bitta musbat, bitta manfiy ildiz.\\
2) $x^2-x-\sqrt3=0$: \ $D=1+4\sqrt3>0$, köpaytma $-\sqrt3<0$ $\Rightarrow$ bitta musbat, bitta manfiy ildiz.\\
Ildizlar ustma-ust tuşmaydi (aks holda $x^2+x+1=x^2-x\Rightarrow x=-\frac12$, lekin $\frac14+\frac12=\frac34\ne\sqrt3$).

a) Musbat ildizlar: 2 ta. \qquad b) Haqiqiy ildizlar: 4 ta.
\Javob{a) 2; \quad b) 4}

%================= 37
\newpage
\Savol{37}
Öng tomon: $2\sin\left(\frac\pi4+x\right)\sin\left(\frac\pi4-x\right)=\cos2x-\cos\frac\pi2=\cos2x$.\\
Çap tomon: $\tg2x-\tg x=\dfrac{\sin x}{\cos2x\cos x}$, \ $\tg2x\tg x=\dfrac{\sin2x\sin x}{\cos2x\cos x}$ $\Rightarrow$ nisbat $=\sin2x$ \ ($\sin x\ne0$, $\cos x\ne0$, $\cos2x\ne0$).
\[\sin2x=\cos2x\iff\tg2x=1\iff x=\frac\pi8+\frac{\pi k}2.\]
Bu nuqtalarda $\sin x\ne0$, $\cos x\ne0$, $\cos2x\ne0$ — hammasi yaroqli.

a) Eng katta manfiy: $\frac\pi8-\frac\pi2=-\frac{3\pi}8$.

b) $[0;\pi]$ da: $\frac\pi8,\ \frac{5\pi}8$ — 2 ta.
\Javob{a) $-\frac{3\pi}8$; \quad b) 2}

%================= 38
\newpage
\Savol{38}
$x=1$ va $x=-1$ ni qöyamiz:
\[x=1:\ 2f(0)+0=2f(1)\Rightarrow f(1)=f(0);\qquad x=-1:\ 0-2f(0)=-2f(-1)\Rightarrow f(-1)=f(0).\]
a) $f(1)-f(-1)=0$.

b) $f(1)=f(0)$: $1+a+b=0$; \ $f(-1)=f(0)$: $-1+a-b=0$ $\Rightarrow$ $a=0$, $b=-1$, \ $f(x)=x^3-x+c$.\\
(Tekşiriş: $(x+1)f(x-1)+(x-1)f(x+1)=2x^4-2x^2+2cx=2xf(x)$ — töğri.)
\[f(2)-f(-2)=(8-2+c)-(-8+2+c)=12.\]
\Javob{a) 0; \quad b) 12}

%================= 39
\newpage
\Savol{39}
Grafikdan: $f'(x)$ nollari va işora özgarişi: $x=-5$ ($-\to+$), $x=-1$ ($+\to-$), $x=1$ ($-\to+$), $x=3$ ($+\to-$). $x>3$ da $f'<0$ (nolga yetmaydi).

a) $f$ ning ekstremum nuqtalari — $f'$ işora almaştiradigan nuqtalar: 4 ta.

b) $\dfrac{f'}{f''}<0$ $\iff$ $f'$ va $f''$ har xil işorali, ya'ni $f'\ne0$ va $|f'|$ kamayadi (grafik $Ox$ ga yaqinlaşadi).\\
Butun nuqtalarni tekşiramiz:
\begin{itemize}\setlength{\itemsep}{0pt}
\item $x=-6,-3,0,2,4,5,6,7,8$ — $f'$ ning ekstremumlari, $f''=0$ — yaroqsiz;
\item $x=-5,-1,1,3$ — $f'=0$, nisbat $0$ — yaroqsiz;
\item $x=-4$: $f'>0$, ösadi ($f''>0$) — yaroqsiz;
\item $x=-2$: $f'>0$, kamayadi ($f''<0$) — yaroqli.
\end{itemize}
Faqat $x=-2$ — 1 ta.
\Javob{a) 4; \quad b) 1}

%================= 40
\newpage
\Savol{40}
Soha: $-1\le x\le0$, \ $-1\le y\le x^3$ ($y=x^3$, $x=0$ va $y=-1$ bilan çegaralangan).
\fig{\begin{tikzpicture}[scale=2]
\fill[okgreen!30,domain=-1:0,samples=40] (-1,-1) -- plot (\x,{\x*\x*\x}) -- (0,-1) -- cycle;
\draw[->] (-1.4,0)--(0.6,0) node[below]{$x$}; \draw[->] (0,-1.3)--(0,0.6) node[left]{$y$};
\draw[thick,domain=-1.08:0.8,samples=50] plot (\x,{\x*\x*\x}) node[right]{$y=x^3$};
\draw (-1.3,-1)--(0.4,-1) node[right]{$y=-1$};
\node[above] at (-1,0){$-1$}; \draw[ax] (-1,0)--(-1,-1);
\end{tikzpicture}}
a) $S=\displaystyle\int_{-1}^0\left(x^3+1\right)dx=\left(\frac{x^4}4+x\right)\Big|_{-1}^0=-\left(\frac14-1\right)=\frac34$.

b) $Oy$ atrofida: kesim — radiusi $|x|=|y|^{\frac13}$ bölgan doira, $y\in[-1;0]$:
\[V=\pi\int_{-1}^0|y|^{\frac23}dy=\pi\cdot\frac35\approx\frac{3\cdot3}5=\frac95.\]
\Javob{a) $\frac34$; \quad b) $\frac95$}

%================= 41
\newpage
\Savol{41}
$\angle B=90^\circ$, $BC=a$, $AB=c$, $AC=b$; $BM$ — mediana, $BL$ — bissektrisa, $BH$ — balandlik.
\fig{\begin{tikzpicture}[scale=0.8]
\coordinate (B) at (0,0); \coordinate (A) at (0,3); \coordinate (C) at (6,0);
\coordinate (M) at (3,1.5); \coordinate (L) at (2,2); \coordinate (H) at (1.2,2.4);
\draw (A)--(B)--(C)--cycle; \draw (B)--(M); \draw (B)--(L); \draw[ax] (B)--(H);
\draw (0,0.3)--(0.3,0.3)--(0.3,0);
\node[left] at (A){$A$}; \node[below left] at (B){$B$}; \node[below right] at (C){$C$};
\node[above right] at (M){$M$}; \node[above] at (L){$L$}; \node[above] at (H){$H$};
\end{tikzpicture}}
Töğri burçakli uçburçakda $BM=\frac b2\Rightarrow b=96$. \ Bissektrisa: $BL=\dfrac{\sqrt2\,ac}{a+c}=8\Rightarrow ac=4\sqrt2(a+c)$.\\
$s=a+c$: \ $s^2-2ac=96^2\Rightarrow s^2-8\sqrt2s-9216=0\Rightarrow s=\dfrac{8\sqrt2+136\sqrt2}2=72\sqrt2$, \ $ac=576$.

a) $S_{ABC}=\frac{ac}2=288$.

b) $(a-c)^2=s^2-4ac=10368-2304=8064\Rightarrow|a-c|=24\sqrt{14}$. \ $\frac{AL}{LC}=\frac ca$ $\Rightarrow$ $AL=\frac{bc}{a+c}$:
\[ML=\left|\frac b2-\frac{bc}{a+c}\right|=\frac{b\,|a-c|}{2(a+c)}=\frac{96\cdot24\sqrt{14}}{144\sqrt2}=16\sqrt7,\qquad BH=\frac{ac}b=6.\]
\[S_{BML}=\frac12\cdot ML\cdot BH=\frac12\cdot16\sqrt7\cdot6=48\sqrt7.\]
\Javob{a) 288; \quad b) $48\sqrt7$}

%================= 42
\newpage
\Savol{42}
$AB\perp AD$, $AB=5$ (kiçik yon tomon), $BC=3$, $AD=7$. $\triangle BOC\sim\triangle DOA$: \ $\frac{CO}{OA}=\frac{BC}{AD}=\frac37$, \ $\frac{AO}{AC}=\frac7{10}$.
\fig{\begin{tikzpicture}[scale=0.62]
\coordinate (A) at (0,0); \coordinate (B) at (0,5); \coordinate (C) at (3,5); \coordinate (D) at (7,0); \coordinate (O) at (2.1,3.5); \coordinate (K) at (0,3.5);
\fill[okgreen!30] (C)--(O)--(D)--cycle;
\draw (A)--(B)--(C)--(D)--cycle; \draw (A)--(C); \draw (B)--(D);
\draw[axs] (O)--(K); \draw (0,0.3)--(0.3,0.3)--(0.3,0);
\node[below left] at (A){$A$}; \node[above left] at (B){$B$}; \node[above] at (C){$C$}; \node[below right] at (D){$D$};
\node[below] at (O){$O$}; \node[left] at (K){$K$}; \node[above] at (1.5,5){$3$}; \node[below] at (3.5,0){$7$}; \node[left] at (0,1.7){$5$};
\end{tikzpicture}}
a) $OK\perp AB$, $OK\parallel BC$: \ $\triangle AKO\sim\triangle ABC$ $\Rightarrow$ $OK=BC\cdot\dfrac{AO}{AC}=3\cdot\dfrac7{10}=\dfrac{21}{10}$.

b) $S_{ACD}=\frac12\cdot7\cdot5=\frac{35}2$; \ $S_{COD}=S_{ACD}\cdot\dfrac{CO}{CA}=\dfrac{35}2\cdot\dfrac3{10}=\dfrac{21}4$.
\Javob{a) $\frac{21}{10}$; \quad b) $\frac{21}4$}

%================= 43
\newpage
\Savol{43}
Beşburçak tomoni $s$, içki burçagi $108^\circ$. $\triangle BCD$ teng yonli $\Rightarrow$ $\angle CDB=\angle CBD=36^\circ$, $\angle ABD=72^\circ$. Rasmdan $F$ — $CD$ örtasi, $AF$ — simmetriya öqi, $AF\perp CD$.
\fig{\begin{tikzpicture}[scale=1.35]
\coordinate (C) at (0,1.701); \coordinate (D) at (1.618,0.526); \coordinate (E) at (1,-1.376); \coordinate (A) at (-1,-1.376); \coordinate (B) at (-1.618,0.526);
\coordinate (F) at (0.809,1.1135); \coordinate (P) at (0.382,0.526);
\fill[okgreen!30] (A)--(B)--(P)--cycle; \fill[okgreen!15] (P)--(F)--(D)--cycle;
\draw (A)--(B)--(C)--(D)--(E)--cycle; \draw (B)--(D); \draw (A)--(F);

\node[below left] at (A){$A$}; \node[left] at (B){$B$}; \node[above] at (C){$C$}; \node[right] at (D){$D$}; \node[below right] at (E){$E$};
\node[above right] at (F){$F$}; \node[below right] at (P){$P$};
\node[font=\scriptsize] at (-0.95,-0.05){$S_2$}; \node[font=\scriptsize] at (0.95,0.75){$S_1$};
\node[font=\scriptsize,aux] at (1.2,0.62){$36^\circ$};
\end{tikzpicture}}
$\triangle PFD$: $\angle F=90^\circ$, $\angle D=36^\circ$, $FD=\frac s2$ $\Rightarrow$ $PD=\dfrac{s}{2\cos36^\circ}$. \ $\cos36^\circ=\frac{1+\sqrt5}4$ $\Rightarrow$ $PD=\dfrac{2s}{1+\sqrt5}=\dfrac{(\sqrt5-1)s}2$.\\
Diagonal: $BD=2s\cos36^\circ=\frac{1+\sqrt5}2s$ $\Rightarrow$ $BP=BD-PD=s$.

a) $\frac{(\sqrt5-1)s}2=\sqrt5-1\Rightarrow s=2$, \ $BP=2$.

b) $S_1=\frac12\cdot FD\cdot PF=\frac12\cdot\frac s2\cdot\frac s2\tg36^\circ=\frac{s^2\tg36^\circ}8$, \quad $S_2=\frac12\cdot AB\cdot BP\sin72^\circ=\frac{s^2\sin72^\circ}2$.
\[\frac{S_2}{S_1}=\frac{4\sin72^\circ}{\tg36^\circ}=\frac{8\sin36^\circ\cos36^\circ\cos36^\circ}{\sin36^\circ}=8\cos^236^\circ=\frac{(1+\sqrt5)^2}2=3+\sqrt5.\]
\[S_2=(3-\sqrt5)(3+\sqrt5)=4.\]
\Javob{a) 2; \quad b) 4}

%================= 44
\newpage
\Savol{44}
Şar $C$ da tekislikka urinadi, $AC=20$ — diametr; rasmdan $D$ nuqta $AB$ da yotadi. Konus öqi $B$ dan ötadi, uning $SK$ yasovçisi şarga $D$ da urinadi ($S$ — konus uçi, $K$ — asos aylanasida).
\fig{\begin{tikzpicture}[scale=1.15]
\begin{scope}
\clip (-1.2,-0.6) rectangle (8.2,3.4);
\draw (0,10) circle (10);
\draw (0,20)--(5,0);
\draw[ax] (80/17,20/17)--(0,10);
\end{scope}
\draw (-1.2,0)--(8.2,0);
\coordinate (C) at (0,0); \coordinate (B) at (5,0); \coordinate (D) at (4.706,1.176); \coordinate (K) at (2.5,0); \coordinate (S) at (5,1.333);
\draw[thick] (K)--(S)--(7.5,0);
\draw[ax] (S)--(B); \draw[axs] (C)--(D);
\draw[thick] (5,0) ellipse (2.5 and 0.25);
\fill (D) circle (1pt);
\node[below] at (C){$C$}; \node[below right] at (B){$B$}; \node[above left] at (D){$D$}; \node[below] at (K){$K$}; \node[above] at (S){$S$};
\node[above left,font=\scriptsize] at (1.5,0){$\frac52$};
\end{tikzpicture}}
a) $AC$ diametr $\Rightarrow$ $\angle ADC=90^\circ$, ya'ni $\triangle CDB$ — töğri burçakli ($\angle D=90^\circ$).\\
$KC$ va $KD$ — $K$ dan şarga urinmalar $\Rightarrow$ $KC=KD$. Töğri burçakli $\triangle CDB$ da $KD=KC$ bölsa, $K$ — gipotenuza örtasi:
\[KB=KC=KD=\frac52\ \Rightarrow\ r=KB=\frac52.\]
b) $\angle KBD=\beta$, \ $\tg\beta=\frac{AC}{CB}=4$. \ $\triangle KDB$ teng yonli ($KD=KB$) $\Rightarrow$ $\angle SKB=180^\circ-2\beta$:
\[\tg\angle SKB=-\tg2\beta=-\frac{2\cdot4}{1-16}=\frac8{15},\qquad h=SB=KB\cdot\frac8{15}=\frac52\cdot\frac8{15}=\frac43.\]
\Javob{a) $\frac52$; \quad b) $\frac43$}

%================= 45
\newpage
\Savol{45}
$AB=h$. Gorizontga nisbatan $AC$ — $45^\circ$, $AD$ — $30^\circ$ ostida pastga tuşadi.
\fig{\begin{tikzpicture}[scale=0.7]
\coordinate (B) at (0,0); \coordinate (A) at (0,4.05); \coordinate (C) at (4.05,0); \coordinate (D) at (7.01,0);
\draw (-1,0)--(7.8,0); \draw (A)--(B); \draw (A)--(C); \draw (A)--(D);
\draw[ax] (A)--(7.5,4.05);
\draw (0,0.3)--(0.3,0.3)--(0.3,0);
\node[left] at (A){$A$}; \node[below] at (B){$B$}; \node[below] at (C){$C$}; \node[below] at (D){$D$};
\node[left] at (0,2){$h$}; \node[below] at (5.5,0){$30$};
\node[font=\scriptsize] at (1.6,3.25){$45^\circ$}; \node[font=\scriptsize] at (2.3,3.85){$30^\circ$};
\end{tikzpicture}}
$\triangle ABC$: $\angle ACB=45^\circ\Rightarrow BC=h$; \quad $\triangle ABD$: $\angle ADB=30^\circ\Rightarrow BD=h\sqrt3$, \ $AD=2h$.

a) $CD=h\sqrt3-h=30\Rightarrow h=\dfrac{30}{\sqrt3-1}=15(\sqrt3+1)\approx15\cdot2{,}7=40{,}5$.

b) $AD=2h=30(\sqrt3+1)\approx81$.
\Javob{a) $15(\sqrt3+1)\approx40{,}5$; \quad b) $30(\sqrt3+1)\approx81$}

\ShowJavoblar
\end{document}
